% Generated draft: requires mathematical and source review.
\documentclass[11pt]{article}
\usepackage[margin=25mm]{geometry}
\usepackage{fontspec}
\setmainfont{Latin Modern Roman}
\usepackage{amsmath,amssymb,mathtools,unicode-math}
\setmathfont{Latin Modern Math}
\usepackage{enumitem,needspace,xurl,hyperref}
\setlength{\parindent}{0pt}
\setlength{\parskip}{5pt}
\setlength{\emergencystretch}{3em}
\newcommand{\N}{\mathbb{N}}
\newcommand{\Z}{\mathbb{Z}}
\newcommand{\Q}{\mathbb{Q}}
\newcommand{\R}{\mathbb{R}}
\newcommand{\C}{\mathbb{C}}
\newcommand{\F}{\mathbb{F}}
\newcommand{\A}{\mathbb{A}}
\newcommand{\PP}{\mathbb{P}}
\newcommand{\E}{\mathbb{E}}
\newcommand{\eps}{\varepsilon}
\newcommand{\dd}{\,\mathrm{d}}
\begin{document}
\section*{UnsolvedMath ID: 3900016}
\textbf{Source problem code:} catalogue entry
\subsection{Definitions and mathematical statement}
For an integer \(n\ge3\), let \(\mathcal P_n\) be the class of convex planar \(n\)-gons. In the standard vertex-triangulation interpretation, a triangulation \(T\) of \(P\in\mathcal P_n\) partitions \(P\) into \(n-2\) triangles using noncrossing diagonals between its original vertices, with disjoint triangle interiors and no added vertices. Let
\[
 d(T)=\left|\{\operatorname{area}(\Delta):\Delta\text{ is a triangle of }T\}\right|,
 \qquad
 A(P)=\max_{T\text{ triangulates }P}d(T).
\]
Areas are ordinary unsigned Euclidean areas; equal numerical areas count only once. The largest pointwise guarantee in the source is
\[
 t(n)=\min_{P\in\mathcal P_n} A(P).
\]
Determine \(t(n)\) for every \(n\ge3\). Equivalently, identify the largest integer \(m\) for which every convex \(n\)-gon has some triangulation with at least \(m\) different triangle areas. Sharpness requires a polygon for which no triangulation has more than \(t(n)\) different areas.

For the lattice-vertex version put
\[
 \mathcal P_n^{\mathbb Z}=\{P\in\mathcal P_n:\text{all vertices of }P\text{ lie in }\mathbb Z^2\},
 \qquad
 t_{\mathbb Z}(n)=\min_{P\in\mathcal P_n^{\mathbb Z}} A(P).
\]
Determine this corresponding largest guarantee as well; lattice vertices do not impose an equal-area requirement or permit additional triangulation vertices.

The displayed definitions use the usual nondegenerate convex-polygon convention, with all listed vertices genuine extreme points. The source does not explicitly address weak convexity with collinear boundary vertices or the possibility of Steiner vertices, so those conventions require confirmation if a different class was intended. The research summary says the unrestricted exact function appears open, with nontrivial bounds and lattice-case results known, but supplies neither their values nor references. It does not establish that the lattice function has been completely determined. The source index's open label was not checked item by item against later literature.
\subsection{Short English statement}
Choose a triangulation of each convex \(n\)-gon to maximize the number of different triangle areas, then ask for the worst polygon's optimum. Determine that guarantee both for arbitrary polygons and for polygons whose vertices are lattice points.
\subsection{Sources}
\begin{itemize}
\item Local registry: \url{data/unsolvedmath/problems.json}, permanent UnsolvedMath ID \(3900016\), source code catalogue entry; supplied statement, background, and research summary, updated 2026-08-06. Attribution: David Eppstein and cited proposers.
\item David Eppstein, \emph{The Geometry Junkyard: Open Problems}, ``Triangulations with many distinct areas'': \url{https://ics.uci.edu/~eppstein/junkyard/open.html}. HTML extraction accessed 2026-07-29. Later item-level resolution was not checked.
\item The research summary mentions nontrivial bounds and lattice-vertex special-case results, but supplies no numerical bounds or identifiable original publications.
\end{itemize}
\end{document}
